\documentclass[10pt,a4paper]{amsart}
\usepackage{amsmath,amssymb,amsthm,graphicx}
% Retained preamble for the Stacks Project chapter snapshot in this directory.
% The chapter's own preamble.tex is a copy of the shared Stacks Project preamble; the block
% below is sliced MECHANICALLY, line for line, from it. Nothing was added, removed or
% redefined; the wrapper class, the xr-hyper/hyperref wiring and the externaldocument list
% are omitted because the assembler supplies the document class.
%
% lines taken from preamble.tex: 3, 5, 9, 11, 12, 156, 157, 158, ... (51 lines total)
%
% package coverage restored from the chapter's own preamble (2026-10-13): xy, xy-2cell

\usepackage[all]{xy}
\xyoption{2cell}
\UseAllTwocells

\newenvironment{reference}{\comment}{\endcomment}
\newenvironment{slogan}{\comment}{\endcomment}
\newenvironment{history}{\comment}{\endcomment}
\theoremstyle{plain}
\newtheorem{theorem}[subsection]{Theorem}
\newtheorem{proposition}[subsection]{Proposition}
\newtheorem{lemma}[subsection]{Lemma}
\theoremstyle{definition}
\newtheorem{definition}[subsection]{Definition}
\newtheorem{example}[subsection]{Example}
\newtheorem{exercise}[subsection]{Exercise}
\newtheorem{situation}[subsection]{Situation}
\theoremstyle{remark}
\newtheorem{remark}[subsection]{Remark}
\newtheorem{remarks}[subsection]{Remarks}
\numberwithin{equation}{subsection}
\def\lim{\mathop{\mathrm{lim}}\nolimits}
\def\colim{\mathop{\mathrm{colim}}\nolimits}
\def\Spec{\mathop{\mathrm{Spec}}}
\def\Hom{\mathop{\mathrm{Hom}}\nolimits}
\def\Ext{\mathop{\mathrm{Ext}}\nolimits}
\def\SheafHom{\mathop{\mathcal{H}\!\mathit{om}}\nolimits}
\def\SheafExt{\mathop{\mathcal{E}\!\mathit{xt}}\nolimits}
\def\Sch{\mathit{Sch}}
\def\Mor{\mathop{\mathrm{Mor}}\nolimits}
\def\Ob{\mathop{\mathrm{Ob}}\nolimits}
\def\Sh{\mathop{\mathit{Sh}}\nolimits}
\def\NL{\mathop{N\!L}\nolimits}
\def\CH{\mathop{\mathrm{CH}}\nolimits}
\def\proetale{{pro\text{-}\acute{e}tale}}
\def\etale{{\acute{e}tale}}
\def\QCoh{\mathit{QCoh}}
\def\Ker{\mathop{\mathrm{Ker}}}
\def\Im{\mathop{\mathrm{Im}}}
\def\Coker{\mathop{\mathrm{Coker}}}
\def\Coim{\mathop{\mathrm{Coim}}}
\DeclareMathSymbol{\boxtimes}{\mathbin}{AMSa}{"02}
\def\QCohstack{\mathcal{QC}\!\mathit{oh}}
\def\Cohstack{\mathcal{C}\!\mathit{oh}}
\def\Spacesstack{\mathcal{S}\!\mathit{paces}}
\def\Quotfunctor{\mathrm{Quot}}
\def\Hilbfunctor{\mathrm{Hilb}}
\def\Curvesstack{\mathcal{C}\!\mathit{urves}}
\def\Polarizedstack{\mathcal{P}\!\mathit{olarized}}
\def\Complexesstack{\mathcal{C}\!\mathit{omplexes}}
\def\Pic{\mathop{\mathrm{Pic}}\nolimits}
\def\Picardstack{\mathcal{P}\!\mathit{ic}}
\def\Picardfunctor{\mathrm{Pic}}
\def\Deformationcategory{\mathcal{D}\!\mathit{ef}}
\usepackage{cleveref}
\crefname{originallemma}{original lemma}{original lemmas}
\Crefname{originallemma}{Original lemma}{Original lemmas}
\makeatletter
\newlabel{lemma-finite-flat-trace}{{0BT8}{source}{}{sourcepointer-1}{}}
\newlabel{lemma-finite-flat-trace@cref}{{[originallemma][0][]0BT8}{[1][0][]source}{}{}{}}
\makeatother
\crefname{originallemma}{original lemma}{original lemmas}
\Crefname{originallemma}{Original lemma}{Original lemmas}
\makeatletter
\newlabel{lemma-trace-base-change}{{0BT9}{source}{}{sourcepointer-2}{}}
\newlabel{lemma-trace-base-change@cref}{{[originallemma][0][]0BT9}{[1][0][]source}{}{}{}}
\makeatother
\crefname{originallemma}{original lemma}{original lemmas}
\Crefname{originallemma}{Original lemma}{Original lemmas}
\makeatletter
\newlabel{lemma-trace-unique}{{0BT7}{source}{}{sourcepointer-3}{}}
\newlabel{lemma-trace-unique@cref}{{[originallemma][0][]0BT7}{[1][0][]source}{}{}{}}
\makeatother
\title{Stacks, Discriminants: lemma-glue-trace}
\author{The Stacks Project Authors}
\date{stacks-project:discriminant.tex:local-snapshot-20261012; GNU FDL 1.2 or later}
\begin{document}
\maketitle
\noindent\textit{Editorial note.} Statement span is lines 954-961 of sources/stacks-discriminant/pinned-original-source.tex; proof span is lines 963-995, the complete author proof (31 body lines).\par\medskip
\noindent\textit{Editorial note.} Navigation stubs: 3 cited-but-not-included labels are supplied as non-load-bearing pointers with their pinned tags where the snapshot tag index has them.\par\medskip
\noindent\textit{Editorial note.} extra\_packages is empty on purpose: the retained preamble already loads xy with options and a second declaration clashes in the assembler wrapper.\par\medskip
\noindent\textit{Editorial note.} Text is retained verbatim, including the author's own typos and placeholders.\par\medskip
\noindent\textit{Editorial note.} Queue index 12 of queue\_discriminant.json (pb=1591, stmt 954-961, proof 963-995).\par\medskip
\noindent\textit{Allowed standard prerequisites (editorial).}\begin{itemize}
\item Explicit prerequisite boundary. The cited inputs are prerequisites; the argument in the retained span is not replaced by any cited result.
\end{itemize}\medskip
\ifdefined\hypertarget\hypertarget{sourcepointer-1}{}\fi
\noindent\textit{Source reference.} Original lemma 0BT8: sources/stacks-discriminant/pinned-original-source.tex, source lines 874--877. Cited by the retained statement or proof but not part of the collected spans; supplied only as a navigation pointer. The argument does not depend on its content.\par\medskip
\ifdefined\hypertarget\hypertarget{sourcepointer-2}{}\fi
\noindent\textit{Source reference.} Original lemma 0BT9: sources/stacks-discriminant/pinned-original-source.tex, source lines 894--907. Cited by the retained statement or proof but not part of the collected spans; supplied only as a navigation pointer. The argument does not depend on its content.\par\medskip
\ifdefined\hypertarget\hypertarget{sourcepointer-3}{}\fi
\noindent\textit{Source reference.} Original lemma 0BT7: sources/stacks-discriminant/pinned-original-source.tex, source lines 843--846. Cited by the retained statement or proof but not part of the collected spans; supplied only as a navigation pointer. The argument does not depend on its content.\par\medskip
\begin{lemma}
\label{lemma-glue-trace}
Let $A \to B$ be a flat quasi-finite map of Noetherian rings.
Let $g_1, \ldots, g_m \in B$ be elements generating the unit ideal.
Let $\tau \in \omega_{B/A}$ be an element whose image in
$\omega_{B_{g_i}/A}$ is a trace element for $A \to B_{g_i}$.
Then $\tau$ is a trace element.
\end{lemma}

\begin{proof}
Suppose we have $A \to A_1$ with $A_1$ Noetherian and a product
decomposition $B \otimes_A A_1 = C \times D$ with $A_1 \to C$ finite.
We have to show that the image of $\tau$ in $\omega_{C/A_1}$ is
$\text{Trace}_{C/A_1}$. Observe that $g_1, \ldots, g_m$
generate the unit ideal in $B_1 = B \otimes_A A_1$ and that
$\tau$ maps to a trace element in $\omega_{(B_1)_{g_i}/A_1}$
by Lemma \ref{lemma-trace-base-change}. Hence we may replace
$A$ by $A_1$ and $B$ by $B_1$ to get to the situation as described
in the next paragraph.

\medskip\noindent
Here we assume that $B = C \times D$ with $A \to C$ is finite.
Let $\tau_C$ be the image of $\tau$ in $\omega_{C/A}$.
We have to prove that $\tau_C = \text{Trace}_{C/A}$ in $\omega_{C/A}$.
By the compatibility of trace elements with products
(Lemma \ref{lemma-trace-base-change})
we see that $\tau_C$ maps to a trace element in $\omega_{C_{g_i}/A}$.
Hence, after replacing $B$ by $C$ we may assume that $A \to B$
is finite flat.

\medskip\noindent
Assume $A \to B$ is finite flat. In this case $\text{Trace}_{B/A}$
is a trace element by Lemma \ref{lemma-finite-flat-trace}.
Hence $\text{Trace}_{B/A}$ maps to a trace element in
$\omega_{B_{g_i}/A}$ by Lemma \ref{lemma-trace-base-change}.
Since trace elements are unique (Lemma \ref{lemma-trace-unique})
we find that $\text{Trace}_{B/A}$ and $\tau$ map
to the same elements in $\omega_{B_{g_i}/A} = (\omega_{B/A})_{g_i}$.
As $g_1, \ldots, g_m$ generate the unit ideal of $B$ the map
$\omega_{B/A} \to \prod \omega_{B_{g_i}/A}$ is injective
and we conclude that $\tau_C = \text{Trace}_{B/A}$ as desired.
\end{proof}

\end{document}
